\section{Closing the Hamiltonian model}
\label{sec:hamiltonian}

We now turn a Morse function with eight quadratic variables into a
Hamiltonian map on a closed manifold.  Each of the sixteen products of
north and south sphere caps will contain one copy of its critical set.
The rotations outside the caps must be chosen with care: their endpoint
maps agree with the quadratic model at both poles, although their paths
can differ by full turns. We track this difference to identify the
homotopy classes of the resulting fixed-point loops as well.

\begin{proposition}\label{prop:cap-transfer}
Let $(X,\omega_X)$ be a closed connected symplectic four-manifold. Write
$(u,v)\in\R^4\times\R^4$ and set
\[
 Q(u,v)=\frac{|v|^2-|u|^2}{2},
 \qquad r^2=|u|^2+|v|^2.
\]
Suppose that $F\colon X\times\R^8\to\R$ is a smooth Morse function
such that $F=Q$ whenever $r\ge3$, and let $N=\#\Crit(F)$.
Let $\sigma$ be the standard area form of total area $4\pi$ on the unit
sphere. For every $\lambda>25/2$, the closed symplectic manifold
\[
 (M,\omega_M)=
 \left(X\times(S^2)^4,
 \omega_X+\lambda\sum_{j=1}^4\sigma_j\right)
\]
admits a Hamiltonian $H\in C^\infty((\R/\mathbb Z)\times M,\R)$ such
that every fixed point of $\phi_H^1$ is nondegenerate, every loop
$t\mapsto\phi_H^t(p)$ based at a fixed point $p$ is contractible in
$M$, and
\[
 \#\Fix_0(\phi_H^1;H)=\#\Fix(\phi_H^1)=16N.
\]
Here $\sigma_j$ denotes the pullback of $\sigma$ from the $j$th sphere,
and Hamiltonian vector fields satisfy $\omega_M(X_H,\cdot)=dH$.
\end{proposition}

We first record the local short-period estimate used below. Yorke proved
the sharp lower bound $2\pi/L$ for a nonconstant periodic solution of a
Euclidean vector field with Lipschitz constant $L$~\cite{Yorke1969}.
The weaker coordinate estimate here suffices on a compact invariant set.

\begin{lemma}\label{lem:short-period}
Let $V$ be a smooth vector field on a smooth manifold, with flow
$\Phi^t$, and let $K$ be a compact set on which this flow is defined
for all times and which it preserves. There is a number $\tau_0>0$
such that, for $p\in K$ and $0<\tau<\tau_0$,
\[
 \Phi^\tau(p)=p\quad\Longleftrightarrow\quad V(p)=0.
\]
\end{lemma}

\begin{proof}
Choose finitely many charts with image the Euclidean ball of radius
$3$, whose inverse images of the unit ball cover $K$. Let $C\ge1$
bound both the coordinate vector fields and their first derivatives
on the closed balls of radius $2$, and put $\tau_0=1/(2C)$.
If $0<\tau<\tau_0$, a trajectory
$t\mapsto\Phi^{\tau t}(p)$, $0\le t\le1$, with $p\in K$ stays
in the radius-$2$ ball of a chart whose unit ball contains $p$:
a first exit would require coordinate displacement greater than $1$,
whereas the speed bound gives displacement at most $\tau C<1/2$.
If the trajectory is a loop, write $w$ for its coordinate expression
and $\widehat V$ for the coordinate vector field. Then
\[
 \dot w(t)=\tau\widehat V(w(t)),\qquad
 \int_0^1\widehat V(w(t))\,dt=0.
\]
Set $S=\max_t|\widehat V(w(t))|$. Since
$|w(t)-w(a)|\le\tau S$ for $a,t\in[0,1]$, subtracting the zero
average and using the Lipschitz bound on this convex ball gives
\[
 |\widehat V(w(t))|
 \le\int_0^1|\widehat V(w(t))-\widehat V(w(a))|\,da
 \le C\tau S.
\]
Thus $S=0$, so the loop is constant and $V(p)=0$. The converse follows
from uniqueness of solutions. No homotopy assumption on the loop is used.
\end{proof}

\begin{proof}[Proof of Proposition~\ref{prop:cap-transfer}]
\emph{The local flows.}
Identify $\R^8$ with $\C^4$ by writing
\[
 u=(x_1,y_1,x_2,y_2),\qquad
 v=(x_3,y_3,x_4,y_4),\qquad z_j=x_j+iy_j,
\]
and equip $X\times\C^4$ with
$\Omega=\omega_X+\sum_jdx_j\wedge dy_j$.
Put $s_1=s_2=-1$ and $s_3=s_4=1$, so that
$Q=\frac12\sum_js_j|z_j|^2$.
Let $\Phi_F^\tau$ and $\Phi_Q^\tau$ denote the autonomous Hamiltonian
flows for this form and the stated sign convention.  In a standard
plane, $X_H=(H_y,-H_x)$; consequently
\begin{equation}\label{eq:quadratic-rotation}
 \Phi_Q^\tau(x,z_1,\ldots,z_4)
 =\bigl(x,e^{-is_1\tau}z_1,\ldots,e^{-is_4\tau}z_4\bigr).
\end{equation}

The vector fields of $F$ and $Q$ agree on $r\ge3$, including its
boundary.  That boundary is invariant, so uniqueness of solutions in
both time directions makes $r<3$ invariant as well.  Trajectories there
remain in the compact set $X\times\{r\le3\}$, whereas exterior
trajectories are the explicit rotations in
\eqref{eq:quadratic-rotation}.  Thus both flows are complete.  They
preserve the open set
\[
 U_0=X\times\{z\in\C^4:|z_j|<4\text{ for all }j\}
\]
and its compact closure $K$: points with $r<3$ have every $|z_j|<3$,
and exterior trajectories preserve each $|z_j|$.
Since $dQ\ne0$ on $r\ge3$, all critical points of $F$ lie in $r<3$.
They form a closed discrete subset of $X\times\{r\le3\}$ and hence
are finite.

\smallskip
\noindent\emph{One small time for all fixed points.}
We choose $\epsilon>0$ so that the fixed points of
$\Phi_F^\epsilon$ in $K$ are precisely the critical points of $F$,
all nondegenerate. Lemma~\ref{lem:short-period}, applied to $X_F$ on
the compact invariant set $K$, supplies $\tau_0>0$ for the first
assertion. The zero set of $X_F$ equals $\Crit(F)$ because $\Omega$
is nondegenerate. We next choose the time small enough to obtain
nondegeneracy of the return map as well.

At a critical point $p$, put $A_p=D_pX_F$.  Differentiating
$\Omega(X_F,\cdot)=dF$ at $p$, where $X_F(p)=0$, shows that $A_p$
is the composition of the nonsingular Hessian of $F$ with a symplectic
identification, and is therefore invertible.  Moreover,
\[
 D_p\Phi_F^\epsilon=\exp(\epsilon A_p).
\]
Let $R\ge1$ bound the moduli of all complex eigenvalues of the finitely
many matrices $A_p$.  If $0<\epsilon R<2\pi$, no nonzero eigenvalue
$\mu$ of any $A_p$ satisfies $\epsilon\mu\in2\pi i\mathbb Z$.
The eigenvalues of $\exp(\epsilon A_p)$ therefore all differ from one.
Fix, once and for all,
\begin{equation}\label{eq:small-epsilon}
 0<\epsilon<\min\{\pi,\tau_0,2\pi/R\}.
\end{equation}
This proves the required fixed-point and nondegeneracy assertions for
the local time-$\epsilon$ map.

\smallskip
\noindent\emph{Sphere caps and matching rotations.}
On each sphere use polar coordinates $(\vartheta,\varphi)$ and write
$h=\cos\vartheta$.  Its symplectic form is
\[
 \lambda\sigma
 =\lambda\sin\vartheta\,d\vartheta\wedge d\varphi
 =-\lambda\,dh\wedge d\varphi.
\]
The north and south cap coordinates are, respectively,
\begin{equation}\label{eq:pole-coordinates}
 z_N=\sqrt{2\lambda(1-h)}\,e^{i\varphi},
 \qquad
 z_S=\sqrt{2\lambda(1+h)}\,e^{-i\varphi}.
\end{equation}
These are smooth coordinates at their poles.  Indeed, in the usual
embedding coordinates $(a,b,h)$ of the unit sphere, they are
\[
 z_N=\sqrt{\frac{2\lambda}{1+h}}(a+ib),
 \qquad
 z_S=\sqrt{\frac{2\lambda}{1-h}}(a-ib).
\]
To check the area forms, write $z=R_ze^{i\psi}$.  At the north pole
$R_zdR_z=-\lambda dh$ and $d\psi=d\varphi$; at the south pole
$R_zdR_z=\lambda dh$ and $d\psi=-d\varphi$.  In both cases
$dx\wedge dy=R_zdR_z\wedge d\psi=-\lambda dh\wedge d\varphi$.

The two radius-$5$ caps are disjoint because $\lambda>25/2$:
their height intervals are
\[
 h>1-\frac{25}{2\lambda}
 \quad\text{and}\quad
 h<-1+\frac{25}{2\lambda}.
\]
For each choice $\nu\in\{N,S\}^4$, let $U_\nu\subset M$ be the
product of $X$ with the corresponding four radius-$4$ caps, and let
$\iota_\nu\colon U_0\to U_\nu$ be the symplectic identification
given by \eqref{eq:pole-coordinates}.  The sixteen sets $U_\nu$ are
pairwise disjoint.

For each $j$ choose a smooth function
$\alpha_j\colon[-1,1]\to(0,2\pi)$, constant on the radius-$5$ caps,
with the following values:
\[
\begin{array}{c|cc}
 &\text{north}&\text{south}\\ \hline
 s_j=-1&\epsilon&2\pi-\epsilon\\
 s_j=1&2\pi-\epsilon&\epsilon
\end{array}
\]
Smooth interpolation is possible within
$[\epsilon,2\pi-\epsilon]\subset(0,2\pi)$.
Define the Hamiltonian on $M$
\[
 P=\sum_{j=1}^4P_j(h_j),\qquad
 P_j(h)=\lambda\int_0^h\alpha_j(\xi)\,d\xi,
\]
and denote its flow by $\Psi_t$.  Since
\[
 \iota_{\alpha_j(h)\partial_\varphi}
       (-\lambda\,dh\wedge d\varphi)
 =\lambda\alpha_j(h)\,dh=dP_j,
\]
this flow rotates $\varphi_j$ with positive speed $\alpha_j(h_j)$
and fixes the $X$ coordinate.  A north coordinate rotates by
$e^{i\alpha_jt}$ and a south coordinate by $e^{-i\alpha_jt}$.
The table therefore gives the endpoint multiplier $e^{-is_j\epsilon}$
in either chart, and hence
\begin{equation}\label{eq:cap-endpoint-match}
 \iota_\nu^{-1}\circ\Psi_1\circ\iota_\nu=\Phi_Q^\epsilon
 \quad\text{on }U_0.
\end{equation}
Only the endpoints are equal: for some cap choices the rotation paths
differ by a full turn. Figure~\ref{fig:pole-rotations} illustrates this
distinction in one complex coordinate.

\begin{figure}[ht]
\centering
\begin{tikzpicture}[>=Stealth,every node/.style={font=\small},
                    x=1cm,y=1cm]
\begin{scope}[xshift=-3.1cm]
 \node[align=center] at (0,1.95)
   {North chart: $z_N=R_z e^{i\varphi}$\\
    chart angle $+\epsilon t$};
 \draw[gray!55] (0,0) circle (1.15);
 \draw[gray!65,dashed] (0,0)--(1.15,0);
 \draw[gray!65,dashed] (0,0)--(35:1.15);
 \fill (1.15,0) circle (1.5pt);
 \fill (35:1.15) circle (1.5pt);
 \node[right] at (1.15,0) {$1$};
 \node[above right] at (35:1.15) {$e^{i\epsilon}$};
 \draw[blue!60!black,very thick,->]
   (1.15,0) arc[start angle=0,end angle=35,radius=1.15];
 \node[align=center] at (0,-1.65)
   {$e^{i\epsilon t}$\\short counterclockwise path};
\end{scope}
\begin{scope}[xshift=3.1cm]
 \node[align=center] at (0,1.95)
   {South chart: $z_S=R_z e^{-i\varphi}$\\
    chart angle $-(2\pi-\epsilon)t$};
 \draw[gray!55] (0,0) circle (1.15);
 \draw[gray!65,dashed] (0,0)--(1.15,0);
 \draw[gray!65,dashed] (0,0)--(35:1.15);
 \fill (1.15,0) circle (1.5pt);
 \fill (35:1.15) circle (1.5pt);
 \node[right] at (1.15,0) {$1$};
 \node[above right] at (35:1.15) {$e^{i\epsilon}$};
 \draw[blue!60!black,very thick,->]
   (1.15,0) arc[start angle=0,end angle=-325,radius=1.15];
 \node[align=center] at (0,-1.65)
   {$e^{-i(2\pi-\epsilon)t}=e^{i\epsilon t}e^{-2\pi it}$\\
    long clockwise path};
\end{scope}
\end{tikzpicture}
\caption{The case $s_j=-1$, shown as multiplier paths in the two local
complex charts with angles drawn schematically. Both sphere longitudes increase,
but the south chart reverses their angular orientation. The paths
start at $1$ and end at the same multiplier $e^{i\epsilon}$; the
south path makes one additional negative turn. For $s_j=1$, the
north path makes the additional turn and both endpoints are
$e^{-i\epsilon}$.}
\label{fig:pole-rotations}
\end{figure}

Every $\Psi_t$ preserves all cap radii, so it preserves each $U_\nu$
and the complement of their union.  A nonpolar sphere point cannot be
fixed by $\Psi_1$, because its rotation angle lies strictly between
zero and $2\pi$.  Consequently
\begin{equation}\label{eq:rotation-fixed-set}
 \Fix(\Psi_1)=X\times\{N,S\}^4
 \subset\bigcup_\nu U_\nu.
\end{equation}

\smallskip
\noindent\emph{A correction supported inside the caps.}
We use the elementary composition rule that if Hamiltonian paths
$a_t,b_t$ have generators $A_t,B_t$, then $a_t\circ b_t$ has generator
$A_t+B_t\circ a_t^{-1}$.  Its vector field is
$X_{A_t}+(a_t)_*X_{B_t}$, and symplectic transport gives
$(a_t)_*X_{B_t}=X_{B_t\circ a_t^{-1}}$, proving the rule with our
sign convention.

On the model space define
\begin{equation}\label{eq:local-correction}
 \delta_t=\Phi_F^{\epsilon t}\circ\Phi_Q^{-\epsilon t},
 \qquad
 K_t=\epsilon\bigl(F-Q\circ\Phi_F^{-\epsilon t}\bigr),
 \qquad 0\le t\le1.
\end{equation}
The composition rule shows that $K_t$ generates $\delta_t$.
When $r\ge3$, the backward $F$ flow equals the backward $Q$ flow and
preserves $Q$.  Thus $K_t=0$ and $\delta_t=\operatorname{id}$ there.
Both constituent flows preserve $r<3$, so $\delta_t$ preserves $U_0$.
Uniformly in $t$, the support of $K_t$ is contained in the compact
set $X\times\{r\le3\}$, strictly inside $U_0$.

Define a Hamiltonian $\widetilde K_t$ on $M$ by
\[
 \widetilde K_t|_{U_\nu}=K_t\circ\iota_\nu^{-1},
 \qquad
 \widetilde K_t=0\quad\text{off }\bigcup_\nu U_\nu.
\]
This is smooth jointly in time and space.  In particular it vanishes
on a whole neighborhood of every cap boundary, since a point there
has some $|z_j|$ close to $4$ and therefore $r>3$.
Let $\Delta_t$ be its Hamiltonian flow.  By uniqueness,
\[
 \Delta_t|_{U_\nu}
 =\iota_\nu\circ\delta_t\circ\iota_\nu^{-1},
 \qquad
 \Delta_t=\operatorname{id}\quad\text{off }\bigcup_\nu U_\nu.
\]
Set $\Gamma_t=\Delta_t\circ\Psi_t$.  Combining
\eqref{eq:cap-endpoint-match} and \eqref{eq:local-correction} yields
\begin{equation}\label{eq:inserted-endpoint}
 \iota_\nu^{-1}\circ\Gamma_1\circ\iota_\nu
 =\Phi_F^\epsilon\circ\Phi_Q^{-\epsilon}\circ\Phi_Q^\epsilon
 =\Phi_F^\epsilon.
\end{equation}
On the complement of the cap union, $\Psi_1$ remains in that
complement and $\Delta_1$ acts trivially.  Hence $\Gamma_1=\Psi_1$
there, including at the cap boundaries.  By
\eqref{eq:rotation-fixed-set}, this exterior has no fixed points.
Inside each cap, \eqref{eq:inserted-endpoint} gives precisely the
$N$ critical points of $F$, all lying in $r<3$ and all nondegenerate
by \eqref{eq:small-epsilon}.  Thus every fixed point has been counted,
and $\Gamma_1$ has exactly $16N$ nondegenerate fixed points.

\smallskip
\noindent\emph{The fixed-point loops.}
For each cap choice $\nu$, define the path on $U_0$
\[
 R_{\nu,t}=\Phi_Q^{-\epsilon t}\circ
             \iota_\nu^{-1}\circ\Psi_t\circ\iota_\nu.
\]
The rotation table shows that $R_{\nu,t}$ fixes the $X$ coordinate
and multiplies each $z_j$ by $e^{2\pi i m_{\nu,j}t}$ for an integer
$m_{\nu,j}\in\{-1,0,1\}$. In particular,
$R_{\nu,0}=R_{\nu,1}=\operatorname{id}$, and these rotations preserve
$r$. For a critical point $p=(x,z)$ of $F$, the fixed-point loop of
$\Gamma_t$, expressed in this cap, is
\[
 \iota_\nu^{-1}\Gamma_t\iota_\nu(p)
   =\Phi_F^{\epsilon t}\bigl(R_{\nu,t}(p)\bigr).
\]
Write $R_{\nu,t}(p)=(x,z_\nu(t))$. The fibre loop contracts to $p$
by
\[
 q_s(t)=\bigl(x,(1-s)z_\nu(t)+sz\bigr),
 \qquad 0\le s,t\le1.
\]
Convexity keeps $q_s(t)$ in $r<3$, since $r(p)<3$ and the rotations
preserve $r$. Therefore
\[
 C(s,t)=\Phi_F^{\epsilon t}\bigl(q_s(t)\bigr)
\]
is a homotopy inside $r<3$ from the displayed fixed-point loop to
the constant loop at $p$. It is based at $p$: the endpoint identities
$q_s(0)=q_s(1)=p$ and the fact that $p$ is an equilibrium of $F$
give $C(s,0)=C(s,1)=p$ and $C(1,t)=p$. Thus every fixed-point loop
of $\Gamma_t$ is contractible within its cap.

\smallskip
\noindent\emph{A smooth periodic Hamiltonian.}
The composition rule gives the generator
\[
 L_t=\widetilde K_t+P\circ\Delta_t^{-1}
\]
for $\Gamma_t$.  Choose a smooth nondecreasing map
$\rho\colon[0,1]\to[0,1]$ equal to zero near zero and to one near
one.  The reparametrized path $\Gamma_{\rho(t)}$ is generated by
\[
 H_t=\rho'(t)L_{\rho(t)}
 =\rho'(t)\bigl(\widetilde K_{\rho(t)}
                  +P\circ\Delta_{\rho(t)}^{-1}\bigr).
\]
This Hamiltonian vanishes on neighborhoods of both time endpoints,
so its periodic extension belongs to
$C^\infty((\R/\mathbb Z)\times M,\R)$.  Its time-one map is exactly
$\Gamma_1$, with the same derivatives at all fixed points.
Composing each based contraction with $t\mapsto\rho(t)$ proves
contractibility for the reparametrized fixed-point loops as well.
This proves the proposition.
\end{proof}
